\section{Conclusion}

A hand written FP16 GEMM on the RTX 5070 reaches 89.9 percent of cuBLAS at 4096
cubed and is demonstrably compute bound, arrived at through nine logged diagnostic
rounds with one change each. The supporting families are correct and benchmarked,
the roofline is built on hardware ceilings rather than on a vendor library, and
every number in this document is regenerated from a results file in the tree by a
script that fails rather than reusing a stale picture.

The finding I would keep if I could only keep one is the negative one. Shared
memory tiling, the first thing every course teaches, is 17 percent slower than the
naive kernel on this part, because the cache hierarchy already absorbs the traffic
that tiling exists to remove. That is a single GPU result and Section
\ref{sec:limitations} says so, but it is the clearest illustration of the method:
an optimization is only worth what the bottleneck it removes was costing, and the
profiler is what tells you which bottleneck that is.

What is not finished is stated as plainly as what is. The 1.1.0 mechanisms, the
tile family, the predicated tails, split-K, stream-K, the CUTLASS reference line,
the SpMV matrix suite, FFT and convolution, and reduction and scan, are implemented
and tested and have no throughput numbers, so their tables read pending. The
existing numbers are v1 era and provisional. Closing both of those is one owner
action, the sweep at locked clocks, and the report is wired so that running it
fills the tables in without anybody editing a document.

What the corrections register adds to the method is the step I had skipped:
auditing my own claims against my own data before publishing them.
